```latex
\documentclass{article}
\usepackage{amsmath,amssymb,amsthm}
\usepackage{hyperref}

\newtheorem{theorem}{Theorem}
\newtheorem{lemma}{Lemma}
\newtheorem{proposition}{Proposition}
\theoremstyle{definition}
\newtheorem{construction}{Construction}

\newcommand{\gp}{^{\mathrm{gp}}}
\newcommand{\sh}{\operatorname{Sh}}
\newcommand{\Psh}{\operatorname{PSh}}
\newcommand{\Lan}{\operatorname{Lan}}
\newcommand{\id}{\operatorname{id}}
\newcommand{\bbZ}{\mathbb Z}
\newcommand{\bbN}{\mathbb N}

\title{A counterexample to the comparison statement}
\author{}
\date{}

\begin{document}

\maketitle

\section{Answer}

\noindent
\textbf{No.}  The asserted quasi-isomorphism is not true in general.  We give an explicit counterexample.

We work throughout with monoids written additively.  Recall that for two elements
$x\in M\gp$ and $y\in M$, the notation $x\prec y$ means that there are integers
$m,n$ such that $m y\le x\le n y$ in the translation-invariant order induced by $M$.

\section{Two elementary lemmas}

\begin{lemma}
\label{lem:inf-injective}
Every infinitesimal monoid homomorphism $p:M\to N$ is injective.
\end{lemma}

\begin{proof}
Suppose that $y\in M$ and $p(y)=0$.  Put $x=-y\in M\gp$.  Then $p(x)=0\in N$.
Moreover $x\prec y$: taking $m=-1$ and $n=0$ gives
\[
-y=my\le x=-y\le 0=ny,
\]
because $y\in M$.  Also $p(y)\prec p(x)$, since both are $0$ and $0\prec 0$.
The infinitesimality of $p$ therefore forces $x=-y\in M$.  Since $M$ is sharp,
$y=0$.  Hence $p$ has trivial kernel and is injective.
\end{proof}

\begin{lemma}
\label{lem:key}
Let $p:M\to V$ be an infinitesimal homomorphism to a valuative monoid $V$,
and let $A,B\in M$.
\begin{enumerate}
\item If $p(A)\ge p(B)$, then $2A-B\in M$.
\item If $p(B)\ge p(A)$, then $2B-A\in M$.
\end{enumerate}
\end{lemma}

\begin{proof}
We prove (1); (2) is symmetric.  If $p(A)=0$, then $p(B)=0$, and injectivity
of $p$ gives $A=B=0$, so the claim is trivial.  Assume therefore that
$p(A)>0$.  If $p(B)=0$, then $B=0$ by injectivity, and again the claim is
trivial.  Thus we may assume $p(A)\ge p(B)>0$.

Put $X=2A-B$ and $Y=A+B$.  Since
\[
X+Y=3A\in M,\qquad 2Y-X=3B\in M,
\]
we have $-Y\le X$ and $X\le 2Y$.  Hence $X\prec Y$.

In $V$, we have $p(X)=2p(A)-p(B)$.  Since $p(A)\ge p(B)$,
\[
p(X)=p(A)+\bigl(p(A)-p(B)\bigr)\ge p(A)>0.
\]
Also
\[
p(Y)=p(A)+p(B)\le 2p(A)\le 2p(X).
\]
Thus $p(Y)\prec p(X)$.  By infinitesimality of $p$, we get $X\in M$, i.e.
$2A-B\in M$.
\end{proof}

\section{The counterexample}

Let
\[
P=\bbN^3=\bbN e_1\oplus \bbN e_2\oplus \bbN e_3.
\]
We take $x=e_1$, and we let
\[
b=P[-e_1]^\sharp\cong \bbN^2,
\]
where the quotient kills the face generated by $e_1$.  We write the canonical
map as
\[
\pi:P\longrightarrow b,\qquad
\pi(e_1)=0,\quad \pi(e_2)=\epsilon_2,\quad \pi(e_3)=\epsilon_3,
\]
with $b=\bbN\epsilon_2\oplus \bbN\epsilon_3$.

\subsection{Two finitely presented $P$-monoids}

Define two $P$-monoids as follows.  Let
\[
Q_+=\bigl(P\oplus \bbN d_+\bigr)\big/\!\bigl(d_+ + e_3=2e_2\bigr)
\]
and
\[
Q_-=\bigl(P\oplus \bbN d_-\bigr)\big/\!\bigl(d_- + e_2=2e_3\bigr),
\]
where the quotients are taken in the category of integral, saturated, sharp
monoids.  Both are finitely presented over $P$: each is obtained by adding one
generator and imposing one relation.  Let
\[
f_+:P\longrightarrow Q_+,\qquad f_-:P\longrightarrow Q_-
\]
be the canonical maps.

\subsection{The family $\{f_+,f_-\}$ is a birhizomic cover}

Let $V$ be a valuative monoid and let $v:P\to V$ be a monoid homomorphism.
Put
\[
a=v(e_2),\qquad c=v(e_3).
\]
Since $V$ is valuative, its order is total, so either $a\ge c$ or $c>a$.

\paragraph{Case 1: $a\ge c$.}
Define $g_+:Q_+\to V$ by
\[
g_+(e_i)=v(e_i)\quad(i=1,2,3),\qquad
g_+(d_+)=2a-c.
\]
This is well defined because $a\ge c$ implies $2a-c=a+(a-c)\in V$, and
\[
g_+(d_++e_3)=(2a-c)+c=2a=g_+(2e_2).
\]
Clearly $g_+\circ f_+=v$.

We claim that $f_+$ is birhizomic at $g_+$.  Let
\[
P \xrightarrow{h} M \xrightarrow{p} V
\]
be a factorization of $v$ with $p$ infinitesimal.  By Lemma~\ref{lem:inf-injective},
$p$ is injective.  In particular $p(h(e_2))=a$ and $p(h(e_3))=c$.
Since $a\ge c$, Lemma~\ref{lem:key} gives
\[
2h(e_2)-h(e_3)\in M.
\]
Define $u:Q_+\to M$ by
\[
u(e_i)=h(e_i)\quad(i=1,2,3),\qquad
u(d_+)=2h(e_2)-h(e_3).
\]
Then
\[
u(d_++e_3)=\bigl(2h(e_2)-h(e_3)\bigr)+h(e_3)=2h(e_2)=u(2e_2),
\]
so $u$ is well defined.  Moreover
\[
p\circ u=g_+,\qquad u\circ f_+=h.
\]
If $u'$ is another such map, then $p\circ u'=p\circ u$, and since $p$ is
injective, $u'=u$.  Hence $f_+$ is birhizomic at $g_+$.

\paragraph{Case 2: $c>a$.}
This is symmetric.  Define $g_-:Q_-\to V$ by
\[
g_-(e_i)=v(e_i)\quad(i=1,2,3),\qquad
g_-(d_-)=2c-a.
\]
Then $g_-\circ f_-=v$, and the same argument using Lemma~\ref{lem:key}
shows that $f_-$ is birhizomic at $g_-$.

Thus every valuative point $v:P\to V$ factors through either $Q_+$ or $Q_-$
in the required way.  Therefore
\[
\{f_+:P\to Q_+,\ f_-:P\to Q_-\}
\]
is a birhizomic cover of the terminal object $P$ of $A_P^{\mathrm{op}}$.

\subsection{The cover does not contain the structure map to $b$}

We now show that neither $Q_+$ nor $Q_-$ maps to $b$ over $P$.

Suppose first that $u:Q_+\to b$ is a $P$-monoid homomorphism.  Then
$u(e_2)=\epsilon_2$ and $u(e_3)=\epsilon_3$.  The relation
$d_++e_3=2e_2$ in $Q_+$ gives
\[
u(d_+)+\epsilon_3=2\epsilon_2.
\]
Hence in the group $\bbZ\epsilon_2\oplus \bbZ\epsilon_3$ we have
\[
u(d_+)=2\epsilon_2-\epsilon_3.
\]
But $2\epsilon_2-\epsilon_3\notin \bbN^2=b$, contradiction.  Therefore no such
$u$ exists.

Similarly, if $u:Q_-\to b$ existed, the relation $d_-+e_2=2e_3$ would force
\[
u(d_-)=2\epsilon_3-\epsilon_2\notin b,
\]
again a contradiction.

Let
\[
z=\pi^{\mathrm{op}}:b\longrightarrow P
\]
be the morphism in $A_P^{\mathrm{op}}$ corresponding to $\pi:P\to b$.  The
birhizomic cover $\{f_+,f_-\}$ generates a covering sieve $S$ of the terminal
object $P$.  A morphism $b\to P$ lies in $S$ only if it factors through
$f_+^{\mathrm{op}}$ or $f_-^{\mathrm{op}}$, i.e. only if there is a
$P$-monoid homomorphism $Q_+\to b$ or $Q_-\to b$.  By the previous paragraph,
no such homomorphism exists.  Hence
\[
z\notin S.
\]
Consequently the pullback sieve
\[
R=z^*S
\]
on $b$ is a covering sieve of $b$ that does not contain $\id_b$.

Moreover, every object $U$ appearing in the sieve $R$ has the following
property: there is no morphism $b\to U$ in $A_P^{\mathrm{op}}$.  Indeed, if
$U\in R$, then the composite $U\to b\to P$ lies in $S$, so it factors through
some $Q_\pm$.  Thus there is a morphism $U\to Q_\pm$ in $A_P^{\mathrm{op}}$.
If there were also a morphism $b\to U$, then composing would give a morphism
$b\to Q_\pm$, equivalently a $P$-monoid homomorphism $Q_\pm\to b$, which we
have just shown does not exist.

\section{Sheaf-theoretic consequences}

Let $A=A_P^{\mathrm{op}}$ with the birhizomic topology, and let
$B=B_P^{\mathrm{op}}$ with the trivial topology.  Let
\[
i:B\longrightarrow A
\]
be the inclusion.  Define a presheaf $F$ of abelian groups on $B$ by
\[
F(b)=\bbZ,\qquad F(b')=0\quad\text{for }b'\neq b,
\]
with all nonidentity restriction maps zero.  Since the topology on $B$ is
trivial, $F$ is a sheaf.

Let
\[
G=\Lan_i F
\]
be the left Kan extension of $F$ along $i$.  For any object $U$ of $A$,
\[
G(U)=\varinjlim_{(b'\to U)\in (i\downarrow U)} F(b').
\]
Since $F$ is supported only at $b$, we have $G(U)=0$ whenever there is no
morphism $b\to U$ in $A$.  By the property of the sieve $R$ established above,
\[
G(U)=0\quad\text{for every }U\in R.
\]
Recall that $\pi^{-1}F$ is the sheafification of $G$ on $A$.

\begin{lemma}
\label{lem:vanish}
Let $G$ be a presheaf on $A$, let $b\in B$, and let $R$ be a covering sieve of
$b$ such that $G(U)=0$ for every $U\in R$.  Then the sheafification $aG$
satisfies $aG(b)=0$.
\end{lemma}

\begin{proof}
We use the plus construction.  For a presheaf $H$ and an object $U$,
\[
H^+(U)=\varinjlim_{T\in J(U)} H^0(T,H),
\]
where $J(U)$ is the collection of covering sieves of $U$.

First, if $U\in R$ and $T$ is any covering sieve of $U$, then every object
$V$ in $T$ also lies in $R$, because $R$ is a sieve on $b$.  Since $G$ is zero
on every object of $R$, we get $H^0(T,G)=0$ for every such $T$.  Hence
\[
G^+(U)=0\quad\text{for all }U\in R.
\]
Now use the covering sieve $R$ of $b$.  The covering sieves contained in $R$
are cofinal among all covering sieves of $b$, because the intersection of two
covering sieves is again a covering sieve.  For every covering sieve
$T\subseteq R$, all objects of $T$ lie in $R$, so $G^+(U)=0$ for all $U\in T$.
Thus $H^0(T,G^+)=0$, and therefore
\[
G^{++}(b)=0.
\]
Since the sheafification is the iterated plus construction, $aG(b)=G^{++}(b)=0$.
\end{proof}

Applying Lemma~\ref{lem:vanish} to our $G$ and the covering sieve $R$, we get
\[
\pi^{-1}F(b)=aG(b)=0.
\]
Hence
\[
(\pi_*\pi^{-1}F)(b)=0.
\]
But $F(b)=\bbZ$.  Therefore the unit
\[
F\longrightarrow \pi_*\pi^{-1}F
\]
is not an isomorphism.

Finally, because $F$ and $\pi^{-1}F$ are sheaves concentrated in degree $0$,
the morphism
\[
F\longrightarrow \mathrm R\pi_*\pi^{-1}F
\]
has $H^0$-component at $b$ given by
\[
F(b)=\bbZ \longrightarrow R^0\pi_*(\pi^{-1}F)(b)=\pi_*\pi^{-1}F(b)=0,
\]
which is not an isomorphism.  Thus it is not a quasi-isomorphism.

\medskip

\noindent
This proves that the asserted comparison statement is false in general.

For background on saturated monoids and finite presentation in this category,
see \href{https://arxiv.org/pdf/2411.13662}{arxiv.org}; for related material on
valuativisations and monoid quotients, see
\href{https://pub.math.leidenuniv.nl/~holmesdst/notes_from_talks/valuativisations_talk.pdf}{pub.math.leidenuniv.nl}.

\end{document}
```